\begin{table}[H]
\centering
\caption{Census of Governments municipal debt balance}
\begin{tabular*}{\textwidth}{@{\extracolsep{\fill}}llccc}
\toprule
Year & Variable & Control & Treat & Difference \\
\midrule
2012 & Long-term debt outstanding & 43.33 & 120.53 & 77.20*** \\
 & Total outstanding debt & 46.61 & 121.43 & 74.81*** \\
 & Total outstanding debt per capita & 2,187.06 & 6,395.78 & 4,208.72 \\
 & Population & 20.47 & 39.70 & 19.23*** \\
\addlinespace
2017 & Long-term debt outstanding & 41.30 & 124.93 & 83.63*** \\
 & Total outstanding debt & 44.30 & 125.76 & 81.46*** \\
 & Total outstanding debt per capita & 2,151.45 & 5,915.97 & 3,764.52 \\
 & Population & 20.45 & 41.21 & 20.76*** \\
\addlinespace
2022 & Long-term debt outstanding & 54.06 & 149.22 & 95.16*** \\
 & Total outstanding debt & 56.53 & 150.53 & 94.01*** \\
 & Total outstanding debt per capita & 2,676.68 & 4,419.36 & 1,742.68 \\
 & Population & 20.87 & 42.99 & 22.12*** \\
\midrule
2012 municipalities &  & 976 & 2,043 &  \\
2017 municipalities &  & 977 & 2,045 &  \\
2022 municipalities &  & 977 & 2,046 &  \\
\bottomrule
\end{tabular*}
\begin{minipage}{\textwidth}
\footnotesize Notes: This table reports means for Census of Governments municipal units matched to the Mergent issuer-level sample. Treat equals one when the Mergent issuer is in a state requiring city GO bond referendums. Difference is treated minus control. Stars denote Welch two-sample t-tests: * $p<0.10$, ** $p<0.05$, *** $p<0.01$. Long-term debt outstanding is Census item 49U. Total outstanding debt is end-of-year long-term plus short-term debt, in millions of dollars. Population is in thousands.
\end{minipage}
\end{table}
